% ============================================================================
\section{Method}
\label{sec:method}
% ============================================================================

Given a set of posed images $\{I_i\}_{i=1}^{N}$ of an outdoor scene with known per-image sun directions $\{\mathbf{l}_i\}$ and optional binary sky masks $\{M_i^{\text{sky}}\}$, our goal is to learn a Gaussian scene representation that decomposes geometry, material (albedo), and illumination, enabling relighting under novel sun conditions.

We build on the 2DGS surfel representation~\cite{huang20242dgs} with per-image appearance embeddings from LumiGauss~\cite{joaxkal2024lumigauss}.
Our method replaces the per-image SH environment with an explicit sun lighting model, adds geometry-based shadow computation, and introduces several training improvements.
Figure~\ref{fig:pipeline} provides an overview.

\begin{figure*}[t]
    \centering
    \todo{Pipeline figure: show Gaussians $\rightarrow$ normals $\rightarrow$ N$\cdot$L $\rightarrow$ shadow map $\rightarrow$ composited image}
    \caption{Overview of the \ours{} pipeline. Per-Gaussian normals from the surfel orientation are used for Lambert shading with the sun direction. A shadow map rendered from an orthographic sun camera provides per-Gaussian shadow masks. The shadowed direct light, ambient, and sky SH components are composed and multiplied by learned albedo to produce the final image.}
    \label{fig:pipeline}
\end{figure*}

% ── 3.1 Sun Model ──────────────────────────────────────────────────────────
\subsection{Explicit Directional Sun Model}
\label{sec:sun_model}

We decompose the illumination reaching each Gaussian into three physically-motivated terms:

\begin{equation}
    L(\mathbf{x}, \mathbf{n}) = \underbrace{C_{\text{sun}}(\theta) \cdot I(\mathbf{l}_i) \cdot \mathbf{k}(\mathbf{l}_i) \cdot \max(0, \mathbf{n} \cdot \mathbf{l}_i)}_{\text{direct sunlight}} + \underbrace{\mathbf{c}^{\text{amb}}(\mathbf{l}_i)}_{\text{ambient}} + \underbrace{E_{\text{sky}}(\mathbf{n})}_{\text{global sky}}
    \label{eq:lighting}
\end{equation}

\noindent where $\mathbf{n}$ is the surface normal (derived from the surfel quaternion and scale, \S\ref{sec:normals}), $\mathbf{l}_i$ is the sun direction for image $i$, and:

\paragraph{Sun color prior $C_{\text{sun}}(\theta)$.}
We model the spectral shift of sunlight with elevation $\theta$ using an atmospheric Rayleigh scattering approximation.
The raw transmittance per RGB channel is:
\begin{equation}
    \hat{T}_c = \exp\!\left(-\beta_R \cdot \text{AM}(\theta) \cdot \left(\frac{\lambda_c}{650}\right)^{-4} \cdot s\right)
    \label{eq:sun_color}
\end{equation}
where $\text{AM}(\theta) = 1/\sin(\theta)$ is the air mass (clamped to $10 + 2(5 - \theta)$ for $\theta < 5^\circ$ to avoid singularity at the horizon), $\lambda_c \in \{650, 550, 450\}$\,nm for R, G, B, $\beta_R = 0.0596$ is the Rayleigh extinction coefficient, and $s = 0.15$ is a scattering scale.
To ensure the prior yields neutral white light at zenith, we normalise by the zenith ($\theta = 90^\circ$) transmittance:
\begin{equation}
    C_{\text{sun},c}(\theta) = \operatorname{clamp}\!\left(\frac{\hat{T}_c(\theta)}{\hat{T}_c(90^\circ)},\; 0.1,\; 1.5\right)
\end{equation}
This produces physically plausible warm tones at low elevations and neutral white at zenith.

\paragraph{Scene-global SH lighting parameters.}
Rather than storing per-image scalars---which cannot generalise to unseen sun positions---we represent the three lighting quantities as low-order spherical-harmonic functions of the \emph{sun direction} $\mathbf{l}$:
\begin{align}
    I(\mathbf{l})                       & = \sum_{l=0}^{d}\sum_{m=-l}^{l} a_{lm}\, Y_l^m(\hat{\mathbf{l}}), \label{eq:scene_sh_int}          \\
    \mathbf{k}(\mathbf{l})              & = \sum_{l=0}^{d}\sum_{m=-l}^{l} \mathbf{b}_{lm}\, Y_l^m(\hat{\mathbf{l}}), \label{eq:scene_sh_col} \\
    \mathbf{c}^{\text{amb}}(\mathbf{l}) & = \sum_{l=0}^{d}\sum_{m=-l}^{l} \mathbf{d}_{lm}\, Y_l^m(\hat{\mathbf{l}}), \label{eq:scene_sh_amb}
\end{align}
where $d = 2$ (9 SH coefficients each) by default, giving:
\begin{itemize}
    \item $I(\mathbf{l}) \in \mathbb{R}^+$: sun intensity multiplier (DC initialized to $4.0 / c_0$, $c_0 = 0.282095$)
    \item $\mathbf{k}(\mathbf{l}) \in [0.5, 2.0]^3$: colour correction on $C_{\text{sun}}$, accommodating cloud cover and atmospheric variation (DC initialized to $1.0 / c_0$)
    \item $\mathbf{c}^{\text{amb}}(\mathbf{l}) \in \mathbb{R}^{3}_+$: ambient/indirect illumination (DC initialized to $0.15 / c_0$)
\end{itemize}
Higher-order SH terms start at zero.
The effective sun colour is $C_{\text{sun}}(\theta) \cdot \mathbf{k}(\mathbf{l}) \cdot I(\mathbf{l})$.
This reduces the total lighting parameters to a fixed 39 scalars (9 for intensity, $3{\times}9$ for colour correction and ambient), enforces physical consistency across views with similar sun positions, and---crucially---enables evaluation at novel sun directions unseen during training.

\paragraph{Global sky SH.}
We represent the directional sky illumination as degree-1 spherical harmonics $E_{\text{sky}} \in \mathbb{R}^{3 \times 4}$, shared across all images.
This captures the dominant sky gradient (bright horizon, dark zenith) while remaining low-dimensional enough to constrain the optimization.
Evaluated at normal $\mathbf{n}$:
\begin{equation}
    E_{\text{sky}}(\mathbf{n}) = \sum_{l=0}^{1}\sum_{m=-l}^{l} c_{lm} \cdot Y_l^m(\mathbf{n})
    \label{eq:sky_sh}
\end{equation}

\paragraph{Final per-Gaussian color.}
The HDR intensity is gamma-corrected and multiplied by the learned albedo $\mathbf{a}$:
\begin{equation}
    \text{RGB} = \mathbf{a} \cdot \left(L(\mathbf{x}, \mathbf{n})\right)^{1/\gamma}, \quad \gamma = 2.2
    \label{eq:final_color}
\end{equation}

% ── 3.2 Normals ───────────────────────────────────────────────────────────
\subsection{Surface Normals from Surfel Geometry}
\label{sec:normals}

Each 2DGS surfel is parameterized by a quaternion $\mathbf{q}$ and anisotropic scales $\mathbf{s} \in \mathbb{R}^3$.
The surfel plane normal is the third column of the rotation matrix $R(\mathbf{q})$:
\begin{equation}
    \mathbf{n}_w = R(\mathbf{q})_{:,3}
\end{equation}
We apply a view-dependent sign flip to ensure normals face the camera:
\begin{equation}
    \mathbf{n} = \text{sign}(-\mathbf{n}_w \cdot \mathbf{p}_{\text{view}}) \cdot \mathbf{n}_w
\end{equation}
where $\mathbf{p}_{\text{view}}$ is the Gaussian position in camera coordinates.
This yields consistent camera-facing normals for Lambert shading.

% ── 3.3 Shadow Computation ─────────────────────────────────────────────────
\subsection{Shadow Computation}
\label{sec:shadows}

Shadows are applied exclusively to the direct sunlight term; ambient and sky illumination remain unaffected.
The shadowed intensity for Gaussian $j$ in image $i$ becomes:
\begin{equation}
    L_{\text{shad}} = C_{\text{sun}} \cdot I(\mathbf{l}_i) \cdot \mathbf{k}(\mathbf{l}_i) \cdot (\mathbf{n}\cdot\mathbf{l}_i) \cdot \sigma_j + \mathbf{c}^{\text{amb}}(\mathbf{l}_i) + E_{\text{sky}}(\mathbf{n})
    \label{eq:shadowed}
\end{equation}
where $\sigma_j \in [0, 1]$ is the shadow mask (1 = lit, 0 = shadowed).

We implement shadow computation using shadow mapping.

\subsubsection{Shadow Mapping}
\label{sec:shadow_map}

Our primary method adapts classical shadow mapping~\cite{williams1978casting} to Gaussian splatting.
We construct an orthographic sun camera positioned along $\mathbf{l}_i$ at distance $3\times r_{\text{scene}}$ from the scene center, with frustum half-width $1.5\times r_{\text{scene}}$, where $r_{\text{scene}}$ is the maximum distance from any Gaussian to the scene centroid.

We render this view using the surfel rasterizer to obtain a depth map $D_{\text{sun}} \in \mathbb{R}^{R \times R}$ and alpha map $\alpha_{\text{sun}}$.
Sky Gaussians are excluded by zeroing their opacity using the learned $w_j^{\text{sky}}$ parameter (\S\ref{sec:sky_mask}).

For each Gaussian $j$ with world position $\mathbf{x}_j$, we project into the sun camera's clip space, bilinearly sample the shadow depth map, and compare:
\begin{equation}
    \sigma_j = \begin{cases}
        0 & \text{if } d_j > D_{\text{sun}}(\mathbf{u}_j) + b \;\wedge\; \alpha_{\text{sun}}(\mathbf{u}_j) > 0.1 \\
        1 & \text{otherwise}
    \end{cases}
\end{equation}
where $d_j$ is the projected depth, $\mathbf{u}_j$ is the UV coordinate, and $b$ is a bias term to prevent self-shadowing.
Gaussians outside the shadow frustum default to lit.

% ── 3.3b Sun Direction Calibration ─────────────────────────────────────────
\subsection{Learnable Sun Direction Calibration}
\label{sec:sun_cal}

Sun directions from metadata (e.g.\ solar position databases or SfM-derived estimates) may be noisy or biased.
We optionally refine them during training by attaching a learnable offset $\Delta\mathbf{l}_i \in \mathbb{R}^3$ (initialized to zero) to each camera:
\begin{equation}
    \hat{\mathbf{l}}_i = \frac{\mathbf{l}_i + \Delta\mathbf{l}_i}{\|\mathbf{l}_i + \Delta\mathbf{l}_i\|}
    \label{eq:sun_cal}
\end{equation}
The refined direction $\hat{\mathbf{l}}_i$ replaces $\mathbf{l}_i$ in all downstream computations (Lambert shading, shadow mapping, scene-SH evaluation).
To prevent unconstrained drift, we apply an L2 regulariser:
\begin{equation}
    \mathcal{L}_{\text{sun-cal}} = \lambda_{\text{sc}} \|\Delta\mathbf{l}_i\|^2
    \label{eq:sun_cal_reg}
\end{equation}
Calibration is active only in the window $[T_{\text{sc,start}},\, T_{\text{sc,end}}]$ (default 10\,k--30\,k iterations), giving the geometry time to stabilise before the sun direction is adjusted, and freezing the directions before the final refinement iterations.

% ── 3.4 Learnable Sky Mask ──────────────────────────────────────────────────
\subsection{Learnable Sky-Scene Distinction}
\label{sec:sky_mask}

Each Gaussian carries a continuous parameter $w_j^{\text{sky}} \in [0, 1]$, implemented as a clamped learnable scalar.
This parameter serves a dual role:
\begin{itemize}
    \item \textbf{Shadow casting}: In the shadow map, each Gaussian's opacity is multiplied by $w_j^{\text{sky}}$, so sky Gaussians ($w_j^{\text{sky}} \approx 0$) become transparent to the sun.
    \item \textbf{Shadow receiving}: Sky Gaussians are forced to fully lit ($\sigma_j = 1$) regardless of the shadow computation, preventing the sky from being darkened by shadows.
\end{itemize}

\paragraph{Supervision.}
Given binary sky masks $M^{\text{sky}}$ (1 = non-sky, 0 = sky), we render each view using $w_j^{\text{sky}}$ as the per-Gaussian color and supervise with an L1 loss:
\begin{equation}
    \mathcal{L}_{\text{sky}} = \lambda_{\text{sky}} \left\| \text{Render}(w^{\text{sky}}) - M^{\text{sky}} \right\|_1
    \label{eq:sky_loss}
\end{equation}
This is differentiable through the rasterizer and jointly optimized with all other parameters.

\paragraph{Sky distance regularisation.}
To prevent sky Gaussians from drifting into the scene interior, we add a penalty pushing them away from the scene centre:
\begin{equation}
    \mathcal{L}_{\text{sky-dist}} = \lambda_{\text{sd}} \frac{1}{|\mathcal{S}|}\sum_{j \in \mathcal{S}} \left[\max(0,\; f \cdot r_{\text{scene}} - \|\mathbf{x}_j\|)\right]^2
    \label{eq:sky_dist}
\end{equation}
where $\mathcal{S} = \{j : w_j^{\text{sky}} < 0.5\}$ is the set of identified sky Gaussians, $r_{\text{scene}}$ is the cameras extent, and $f = 4$ (default) is the minimum-distance factor.